% GENERATED FILE. Edit canonical lessons, then run npm run generate:books.
% canonical-source-hash: fnv1a64:bffd20e42624850b
% canonical-lessons: PT-C26-olho, PT-C26-orelha, PT-C26-boca, PT-C26-coracao

\chapter{The Body, Asking How You Are}
\label{ch:pt-body-how-are-you}
\hlchaptermodality{26}

\begin{chapteropening}
\textbf{By the end of this chapter:} \emph{I can name the eye, ear, mouth and heart in Portuguese with the right gender, and say what happened to each Latin root -- kept whole, folded by a sound law, or grown with a suffix.}

The chapter's last lesson closes the set with a one-table review of what happened to four body-part roots (kept, eroded, folded, grown), contrasts coracao's built nasal -ao against pao's and irmao's worn-down one, and reaches back to Chapter 2's Tudo bem/formal practice to give the wellbeing check-in something to answer with.
\end{chapteropening}

\section[o olho]{o olho — the sound law that already gave you a colour}

\label{lesson:PT-C26-olho}



\begin{quote}
New chapter: the body, extended past Chapter 17's \emph{cabeça} and \emph{mão}. The first word runs on a sound change you've already met — you just didn't know its name yet.
\end{quote}

\subsection*{You'll want to know: The word}
\begin{quote}
\textbf{o olho} — the eye. \textbf{Masculine.} \textbf{Dói-me o olho.} — \textquotedblleft{}My eye hurts.\textquotedblright{} (Same \emph{dói-me} pattern as Chapter 17's \emph{cabeça}.)
\end{quote}

The \textbf{-lh-} is the palatal \emph{ly} glide you already own from \emph{trabalhar}, \emph{vinho}, and now \emph{mulher}.

\begin{cousinweb}
\textbf{olho} $\leftarrow$ Latin \emph{\textbf{oculus}} (\textquotedblleft{}eye\textquotedblright{}). Spoken Latin dropped the unstressed middle vowel first — \emph{oculus $\to$ oclu} — leaving a \emph{-c-} jammed right against an \emph{-l-}. That exact cluster, \textbf{-c'l-}, is what regularly folds into Portuguese \textbf{-lh-}: \emph{oclu $\to$ olho}.

You have already seen this happen once, without the mechanism named: Chapter 13's \textbf{vermelho} (\textquotedblleft{}red\textquotedblright{}) comes from \emph{vermiculus}, which loses its own unstressed vowel the same way — \emph{vermiculu $\to$ vermelho} — the identical \textbf{-c'l- $\to$ -lh-} change, just on a different word. \emph{Olho} and \emph{vermelho} rhyme in Portuguese for a real historical reason, not by accident.

English kept the Latin shape whole, in the \textbf{learned} vocabulary: \textbf{ocular}, \textbf{oculist}. As with \emph{ovo}/\emph{ovum} two chapters back, the everyday Portuguese word and the everyday English word parted ways; only English's technical register still shows the root.
\end{cousinweb}

\subsection*{Your turn}
Say these aloud:

\begin{itemize}
\raggedright
  \item \textquotedblleft{}o olho\textquotedblright{} — the same \emph{-lh-} as \emph{trabalhar}, \emph{vinho}, \emph{mulher}
  \item \textquotedblleft{}Dói-me o olho\textquotedblright{} — reusing Chapter 17's \emph{cabeça} pattern
  \item \textquotedblleft{}oculus $\to$ oclu $\to$ olho\textquotedblright{} — the same change as vermiculus $\to$ vermelho
  \item the body words so far — \textquotedblleft{}a cabeça, a mão, o olho\textquotedblright{} — feminine, feminine, masculine
\end{itemize}

\begin{tcolorbox}[breakable,colback=teal!4,colframe=teal!35!black,title={Before you move on}]
Say \textquotedblleft{}the eye,\textquotedblright{} and \textquotedblleft{}my eye hurts.\textquotedblright{} (\textbf{O olho}; \textbf{dói-me o olho}.) What Latin word is it from, and what sound law turns its \emph{-c'l-} cluster into \emph{-lh-}? (\emph{\textbf{Oculus}} — the same \textbf{-c'l- $\to$ -lh-} change as Chapter 13's \emph{vermiculus $\to$ vermelho}.) Is \emph{a cabeça} feminine or masculine, and how do you spell the \emph{s}-sound in it? (\textbf{Feminine}; the \textbf{cedilla ç}.) Why is \emph{a mão} feminine even though it doesn't end in \emph{-a}? (Latin \emph{manus} was \textbf{fourth-declension feminine}.) Which Chapter 25 word named \textquotedblleft{}human being\textquotedblright{} before Portuguese narrowed it to \textquotedblleft{}man\textquotedblright{}? (\textbf{Homem}, $\leftarrow$ \emph{homō}.) Next: a second body word built the exact same way.
\end{tcolorbox}
\section[a orelha]{a orelha — the same fold, a second time}

\label{lesson:PT-C26-orelha}



\begin{quote}
The eye's sound law strikes again, on a different word — and this one also reruns a split you already know from a verb, not a body part.
\end{quote}

\subsection*{You'll want to know: The word}
\begin{quote}
\textbf{a orelha} — the (outer) ear. \textbf{Feminine.} \textbf{Dói-me a orelha.} — \textquotedblleft{}My ear hurts.\textquotedblright{}
\end{quote}

A note on scope: \textbf{orelha} names the visible, outer ear. For hearing itself, or an earache felt inside, Portuguese reaches for a different word, \emph{ouvido} — not needed yet, but worth knowing the visible/inner split exists.

\begin{cousinweb}
\textbf{orelha} $\leftarrow$ Latin \emph{\textbf{auricula}}, literally \textquotedblleft{}\textbf{little ear}\textquotedblright{} — a diminutive of \emph{auris}, \textquotedblleft{}ear.\textquotedblright{} Vulgar Latin dropped its own unstressed vowel — \emph{auricula $\to$ }oricla\emph{ — leaving a \textbf{-c-} jammed against an \textbf{-l-}, the very cluster that folds into Portuguese \textbf{-lh-}, exactly as }oculus\emph{ became }olho\emph{ last lesson.}

Spanish took the \textbf{same} Latin word and folded the same cluster a \textbf{different} way: \emph{auricula $\to$ oreja}. That is not a new pattern — it is the \textbf{identical} split Chapter 5 showed you with \emph{trabalhar} and \emph{trabajar}: one shared consonant cluster, Portuguese resolving it to \textbf{-lh-}, Spanish to \textbf{-j-}, twice now on two completely unrelated words. English's \emph{auricle} and \emph{aural} keep the bare Latin root, the same \textquotedblleft{}everyday word vs. learned word\textquotedblright{} split you've now seen several times.
\end{cousinweb}

\subsection*{Your turn}
Say these aloud:

\begin{itemize}
\raggedright
  \item \textquotedblleft{}a orelha\textquotedblright{} — the same \emph{-lh-} as \emph{olho}
  \item \textquotedblleft{}Dói-me a orelha\textquotedblright{} · \textquotedblleft{}Dói-me o olho\textquotedblright{}
  \item \textquotedblleft{}auricula $\to$ orelha (PT) / oreja (ES)\textquotedblright{} — the trabalhar/trabajar split again
  \item \textquotedblleft{}orelha = outer ear; ouvido = hearing, not taught yet\textquotedblright{}
\end{itemize}

\begin{tcolorbox}[breakable,colback=teal!4,colframe=teal!35!black,title={Before you move on}]
Say \textquotedblleft{}the ear,\textquotedblright{} and \textquotedblleft{}my ear hurts.\textquotedblright{} (\textbf{A orelha}; \textbf{dói-me a orelha}.) What Latin word is it from, and what does it literally mean? (\emph{\textbf{Auricula}}, \textquotedblleft{}\textbf{little ear}.\textquotedblright{}) What Chapter 5 verb pair already showed you this exact same Portuguese-\emph{lh}/Spanish-\emph{j} split? (\textbf{Trabalhar / trabajar}.) Which word does \emph{orelha} NOT cover, and what does that word cover instead? (\textbf{Ouvido} — hearing and the inner ear.) Recall Chapter 25's table: which Latin root did French pick for \textquotedblleft{}woman\textquotedblright{} instead of \emph{mulier}? (\textbf{Fēmina}.) Next: the part of the face this whole conversation comes out of.
\end{tcolorbox}
\section[a boca]{a boca — the cheek that took over the mouth}

\label{lesson:PT-C26-boca}



\begin{quote}
The word for what all this talking has been coming out of — and it started life meaning something narrower.
\end{quote}

\subsection*{You'll want to know: The word}
\begin{quote}
\textbf{a boca} — the mouth. \textbf{Feminine.} \textbf{Dói-me a boca.} — \textquotedblleft{}My mouth hurts.\textquotedblright{}
\end{quote}

\begin{cousinweb}
\textbf{boca} $\leftarrow$ Latin \emph{\textbf{bucca}} — but Classical Latin \emph{bucca} meant specifically \textbf{\textquotedblleft{}cheek,\textquotedblright{}} especially a cheek puffed out, the way you'd blow out air. Latin's actual word for \textquotedblleft{}mouth\textquotedblright{} was \emph{\textbf{ōs, ōris}}.

Everyday spoken Latin generalised \emph{bucca} from \textquotedblleft{}cheek\textquotedblright{} to the whole \textbf{mouth}, and the old word \emph{ōs} survived only in \textbf{learned} vocabulary — English \textbf{oral}, \textbf{orifice}. So Portuguese \emph{boca}, Spanish \emph{boca}, French \emph{bouche} and Italian \emph{bocca} all descend from what was once a word for one puffed-out cheek, while the original \textquotedblleft{}mouth\textquotedblright{} word became a textbook term. English \textbf{buccal} (a dentist's word for the cheek side of a tooth) is the one place the original, narrower sense survives.

If someone asks \emph{Tudo bem?} and you're only \textbf{mais ou menos} — Chapter 2's shrug — the reason might be sitting right here: \emph{Dói-me a boca} — \textquotedblleft{}my mouth hurts\textquotedblright{} — is exactly the kind of thing that turns a \textquotedblleft{}so-so\textquotedblright{} into a \textquotedblleft{}so-so.\textquotedblright{}
\end{cousinweb}

\subsection*{Your turn}
Say these aloud:

\begin{itemize}
\raggedright
  \item \textquotedblleft{}a boca\textquotedblright{} · \textquotedblleft{}dói-me a boca\textquotedblright{}
  \item \textquotedblleft{}bucca — cheek — promoted to mouth; ōs — the old word, now only 'oral'\textquotedblright{}
  \item \textquotedblleft{}Tudo bem? — Mais ou menos, dói-me a boca\textquotedblright{} — Chapter 2's shrug, explained
  \item the ear/mouth pair — \textquotedblleft{}a orelha, a boca\textquotedblright{} — both feminine
\end{itemize}

\begin{tcolorbox}[breakable,colback=teal!4,colframe=teal!35!black,title={Before you move on}]
Say \textquotedblleft{}the mouth,\textquotedblright{} and \textquotedblleft{}my mouth hurts.\textquotedblright{} (\textbf{A boca}; \textbf{dói-me a boca}.) What did Latin \emph{bucca} originally mean, and what was the actual Latin word for \textquotedblleft{}mouth\textquotedblright{}? (\textbf{\textquotedblleft{}Cheek\textquotedblright{}}; the real word was \emph{\textbf{ōs, ōris}}, which survives only in \textbf{oral}.) Give the Chapter 2 shrug, then a reason for it using this lesson's word. (\textbf{Mais ou menos} — e.g. \emph{dói-me a boca}.) Next: the last body word, and where its ending really comes from.
\end{tcolorbox}
\section[o coração]{o coração — a nasal ending that was BUILT, not worn away}

\label{lesson:PT-C26-coracao}



\begin{quote}
The chapter's last word carries the same ending as \emph{pão} and \emph{irmão} — and gets there by the opposite process.
\end{quote}

\subsection*{You'll want to know: The word}
\begin{quote}
\textbf{o coração} — the heart. \textbf{Masculine.} \textbf{Como vai o seu coração?} — a gentle, half-figurative way to ask how someone's doing, alongside Chapter 2's \emph{Tudo bem?}
\end{quote}

\begin{cousinweb}
\textbf{coração} is built on Latin \emph{\textbf{cor, cordis}} (\textquotedblleft{}heart\textquotedblright{}) plus an \textbf{augmentative ending}, from Vulgar Latin \emph{-atiōne} — closer to \textquotedblleft{}\textbf{big heart}\textquotedblright{} than a plain, unadorned one.

Compare Chapter 11's \textbf{pão} and Chapter 10's \textbf{irmão}: both reached their nasal \textbf{-ão} by \textbf{losing} sounds. \emph{Coração} reaches the same ending by the \textbf{opposite} route — a suffix was \textbf{added}, and the suffix itself was nasal. Spanish did the identical building job, landing on \emph{corazón}; French \textbf{coeur} and Italian \textbf{cuore} kept the bare root, no suffix at all.

The bare root \emph{cor, cordis} survives in English \textbf{cordial}, \textbf{courage}, and \textbf{record} — literally \textquotedblleft{}to call back to \textbf{heart}.\textquotedblright{}
\end{cousinweb}

\begin{grammarlens}[title={What you've built: The chapter in one table}]
\noindent
\begin{tabularx}{\linewidth}{@{}>{\raggedright\arraybackslash}X>{\raggedright\arraybackslash}X@{}}
\toprule
\textbf{Portuguese} & \textbf{what happened to the root} \\
\midrule
\textbf{a cabeça} (Ch. 17) & \emph{capitia} kept, whole \\
\textbf{a mão} (Ch. 17) & \emph{manus} lost its \emph{-n-} \\
\textbf{o olho / a orelha} & \emph{-c'l-} folded into \emph{-lh-} \\
\textbf{o coração} & \emph{cor} \textbf{grew} a suffix \\
\bottomrule
\end{tabularx}

Kept, eroded, folded, grown — the same \textquotedblleft{}not one of them agrees\textquotedblright{} shape Chapter 13 found in Portuguese's colours. Chapter 2's \emph{Tudo bem?} and its formal twin \emph{Como está o senhor?} now have somewhere to go: \textbf{Mais ou menos, dói-me a orelha}, or formally, \textbf{Estou bem — e o seu coração?}
\end{grammarlens}

\subsection*{Your turn}
Say these aloud:

\begin{itemize}
\raggedright
  \item \textquotedblleft{}o coração\textquotedblright{} — nasal \emph{-ão}, cedilla \emph{ç}
  \item \textquotedblleft{}Como vai o seu coração?\textquotedblright{}
  \item \textquotedblleft{}pão, irmão: nasal by LOSS; coração: nasal by GROWTH\textquotedblright{}
  \item \textquotedblleft{}cabeça, mão, olho/orelha, coração\textquotedblright{} — kept, eroded, folded, grown
\end{itemize}

\begin{tcolorbox}[breakable,colback=teal!4,colframe=teal!35!black,title={Before you move on}]
Say \textquotedblleft{}the heart.\textquotedblright{} (\textbf{O coração}.) How does its \emph{-ão} differ from \emph{pão}'s? (That one \textbf{lost} a sound; \emph{coração} \textbf{grew} a suffix.) Name one English word keeping the bare root \emph{cor}. (\textbf{Cordial}, \textbf{courage}, or \textbf{record}.) What four things happened to this chapter's four roots? (\textbf{Kept, eroded, folded, grown.}) That closes the body-and-wellbeing set begun in Chapter 2.
\end{tcolorbox}
